\chapterNote{Based on a talk given on November 4, 2021}

{\centering
\chapter{Making the Worst of Circumstances the Best of Practices}
\tocChapterNote{Based on a talk given on November 4, 2021}
\par}

We're in the middle of our Kathina celebration here at Abhayagiri, and Luang Por Pasanno and I have just returned from spending the Vassa at Pacific Hermitage. It's really nice to be back here and to enjoy the camaraderie of bhikkhus, lay friends, supporters, and the broader community. We can rejoice in the flourishing of the good thing that's been happening here, including the many Dhamma reflections that have been offered.

At today's Patimokkha reflection, Luang Por Pasanno brought up the theme of rejoicing in this communal gathering of goodness, the goodness that comes with the coming together of people who are all committed to very high values. He noted that there's so much negativity and cynicism these days in the world around us. We see it in political and social issues. There are many difficulties fueled by the pandemic in the world these days, and people are feeling stressed. Their ways of responding feel natural, but are often marked by negativity and cynicism, which are not useful.

We reflect on the importance of not getting caught in those states of mind. It's so easy to do. It's kind of a tunnel that's really easy to go down. We hear things that excite the mind, and we look for external situations on which to blame the difficulties that we experience. We think that if only we can right that wrong then everything would be okay. It's just a relentless, downhill mind state that doesn't go anywhere if we buy into it.

It's this way with other difficult, obstructive states of mind, too. Sometimes they are unrelenting, and even if we are trying to live a good and useful life, we can still feel bombarded by these mind states. They can be very hard to bear. Whether it's anger, relentless sensual craving, doubts, confusion, delusion, or restlessness, to name just a few, it's so easy to let them expand into an even more negative or cynical mind state. If not projected outward to blame some external source, then they may be projected inwards with self-blaming, a sense of worthlessness, or an inability to make progress on the path.

Sometimes the difficulties, the obstructions that we experience in our practice, just come so hard at us. We get this sense of despair and we wonder, ``What's the point? Why do I put up with this? It's not going to go anywhere.''

We can experience great periods of doubt and difficulty. We can acknowledge that it's a common experience for all of us. Every one of us has gone through periods like this, in one way or another. Knowing this doesn't necessarily make it any easier, but sometimes it's useful to expand beyond the self-referencing and realize that it's not just ``me.''

The question is, of course: How to deal with this? How to use the resources that we have to lift ourselves up out of it, to not stay down in a state of difficulty or despair? What resources can we muster to pull ourselves up and keep going?

I've been reading a very compelling book that somebody gave to me recently. I'd heard of it before but had never taken the time to look into it and read it. It's called \emph{Man's Search for Meaning} by Viktor Frankl, and I've just finished the first section.

Viktor Frankl was a psychiatrist in Austria at the time of World War II. He was among those taken and interned into several concentration camps during the war and was a prisoner there for several years. The book contains his recollections about his state of mind while he was in the concentration camps. Right from the beginning, he was very clear that he wasn't writing this book as an account of all of the atrocities and painful circumstances that occurred while he was there, although, of course, he had to relate some of them to provide the context for his own internal processes. He then expounded on these states of mind and what he saw as the internal mechanisms at work for the other people that he was interned with there.

He was an incredibly compelling and reflective person, and one of the few who was able to transcend the horrors of being in that situation. There were the physical horrors, the extreme deprivation, starvation, torture, extreme cold, hard work, malnutrition, barely enough to keep alive, and all of the diseases that would go around in these extremely crowded conditions.

This was if you were lucky enough to have escaped the extermination chambers, having watched your friends and relatives being taken already. In addition to physical deprivation, there were also the psychological difficulties, the horrors of being psychologically tortured day after day. He was able to stay on top of it enough to honestly experience and allow in what was happening, without moving into a state of denial or repression. He was also determined not to give in and let it break him. He noted a common pattern with many of the prisoners: after the initial horror and fears of being brought into the camps, there was a middle phase of intense indifference and apathy. This state of psychological annihilation, tuning out, seemed to be a way of escaping the emotional pain, just to be able to cope. He could see this happening around him and could see the allure of this coping mechanism, but he did not allow himself to completely collapse into that state of mind.

He had an inner resource, an inner drive to not give in to difficult, horrifying mind states. Over time, he was also able to be of service to his fellow inmates and would offer periods of reflection and support in ways that were under the radar to avoid getting punished and beaten.

It was in these subterranean ways that he was able to help others and to pull himself up time after time, when doubts, confusion, and difficulties arose. It's a very inspiring story that we can all benefit from, reflecting on how he was able to transcend such an extreme situation.

There is this potential to not succumb, to not fall down deeply. Several of the ways that Frankl would do that would be to just not allow thought patterns dwelling on hopelessness and despair, to encourage memories of good things, to allow himself the hope of a future return to freedom and goodness. He would think of a return to more normal states of mind. He would allow his mind to attend to memories that were beneficial and would bring up a sense of hope and warmth, no matter how impossible they seemed. He would direct his thoughts in that direction rather than just succumbing to the present-moment horrors of what was going on, without denying or repressing them, but also not dwelling on them.

Even though he didn't write about it in this way, one of his observations or realizations seemed very Buddhist to me. He described the patterns of stimulus and response that we all operate under and that the key to freedom and happiness is noticing and realizing the moment between stimulus and response. Whether it's coming from physical circumstances, mental, or psychological experiences, there is a gap between stimulus and response. There's a pause, no matter how small, and then the response.

Viktor Frankl's conclusion from his experience was that true freedom is found in being able to navigate that gap, so that we can choose where to go in any particular experience or situation. To me, that insight is the crux of our training, too. Before we respond, we can pause to acknowledge our experience, and then we can choose how we want to direct our attention, how we want to respond. We do have that possibility.

If he can do that in such an extreme situation, I have hope that all of us, myself included, can also do that when obstructions and difficulties come into our lives.

We should take delight in what we've been experiencing the past few days and in general how fortunate we are to have access to this kind of environment that supports the skillful attention towards goodness. Sometimes we can lose our perspective and forget that there are pockets of goodness---monasteries or our homes---where we can create an environment that supports reflection and turning towards the good. The environment supports turning towards the skillful and what is beneficial to ourselves and to other people. Places like this are repositories of skillful means and provide direction on where to turn, where to put our hearts and minds. They nurture the values that we can aspire to, the values that lead to an end of suffering. They help us develop the qualities of virtue, goodness, and generosity. We can cultivate the wisdom and compassion that is required in the difficult circumstances that we experience as human beings. Occasions like coming to the monastery for a period of time help us bring these resources back into our lives, back out into the world. This is the kind of world that we want to create. This is the kind of world that we can create, even in the midst of cynicism and negativity.

So really take refuge in the skillfulness and the good-heartedness of people who are struggling in the same ways that we are, trying to overcome these difficulties that are part of human experience. We all share this aspiration, and there's a set of skills and values that we can nourish and keep to the fore all of the time, with the proper effort. We don't have to succumb.

I've been reflecting on the importance of noticing when we are deeply involved with the perception of self and how this compounds and complicates our world and our experience. In a question-and-answer session with Ajahn Amaro, somebody asked him about the state of boredom. He responded, ``You can't be bored if you are dwelling with the perception of not-self.''

When we are operating through and entangled with the perception of self, we can fall for obstructive mind states, whether it's boredom or one of the many others.

If we can learn to cultivate over time, little by little, the perception of not-self, we have the opportunity to see these obstructive habits and mind states as not who we are. Even though we might still experience them, we don't have to pick them up, own them, and compound them.

From a reflective position, we can understand: ``These are phenomena. These are selfless phenomena happening. It's painful. It's difficult. But I don't need to pick it up as mine.'' When we have that perspective, then we have energy and space to be with them. We have the energy and space to be in the gap between stimulus and response.

It takes the sting out of these obstructive mind states and quells the proliferative tendency. Dwelling in, refreshing, and reflecting on the perception of not-self
is a major antidote to our experience of suffering. There are all the conventional ways of contemplating and developing the perception of not-self. It's also useful to contemplate the commonality of our experience with other people.

Before I was a monk, I remember reading some accounts from Stephen Levine, where he was recounting the experience of a woman who was in the process of dying. She was very ill with a terminal illness, and in addition to the relentless pain and physical discomfort, she was caught in mental resistance as well, with a lot of upset and anger. ``Why is this happening? Why can't I get over this? Why is this so painful?'' Due to the painful circumstances of her condition and the fact that it wasn't going to get any better, she was really suffering.

The struggle and the intense suffering was going on for quite some time. She was reacting to the bodily difficulties and struggling hard with her eventual demise. As this was happening, she was also trying to be reflective and conscious of the experience. At one point, it became so difficult that she went into a kind of transcendent state. She started to have visions and thoughts of all of the other people in the world who were going through similar experiences. She saw in great detail the struggling and suffering of person after person, with their own unique circumstances, time after time, in almost terrifying regularity. Through that experience she had an epiphany: ``This isn't me. This isn't mine. I'm not the only one. This is not \emph{my} suffering. This is \emph{the} suffering.'' The personalization of it as \emph{my} suffering fell away, and she realized and saw the enormity of the situation. This is the condition of the world. It wasn't hers only. It's a universal experience.

Being able to transcend the self-absorption, the self-obsession that this was a unique experience happening to her was her liberating moment, her epiphany. She saw it as \emph{the} suffering.

I've always considered the universal experience to be a really important thing to reflect on for ourselves, too---not just the conventional ways of trying to chip away at this perception of self. We can realize with these experiences, these universal difficulties within the human realm---we're all in it, we're all doing it. It's the nature of the human condition. That opens up a great sense of compassion, and we can let go of the singularity of our own worlds.

Whether it's horrible experiences like being in a concentration camp---the worst we can probably imagine in the human realm---or whether it's the internal struggles of our own minds, we can realize that this is the human condition and that none of us are exempt from these experiences. And it's poignant. It's that very deep understanding and penetration of the First Noble Truth: There is suffering. With stories like Viktor Frankl's and the stories of practitioners---nuns and monks, laywomen and laymen---going through these experiences with honesty, clarity, and openness, there comes the realization that we don't have to succumb. We don't have to buy into it. We can, with effort and the support and encouragement of each other, realize that there's a way beyond this.

There is an ending to this if we want to see it through, if we attend to it in the right way and nourish and nurture the supportive qualities in our life: the goodness that comes from attention to virtue, attention to generosity, attention to sharing and helping each other. We share those abilities as we share the suffering in our lives. And it's this spiritual friendship that is, as the Buddha said, the whole of the holy life, whether it's the holy life in the monastery or the holy life in the householder's world. We have to maintain a refuge in the Sangha, this support. We have to come together. We have to support each other and be there for each other. We can't do it without that support and that refuge.

I just feel so grateful to be part of this community and to see the people coming day after day. If ever I think, ``Am I doing a worthwhile thing? Is there something else that I should be doing? Is there a better way to do this? Am I up to this?'' If I ever have those thoughts, I can reflect on the numbers of people that come here to support us and join in with our practice. People want us to be here, all of us.

People want us to keep doing what we're doing because we do help to keep that alternative alive in the world. We don't have to succumb to the forces of difficulty, negativity, or cynicism. People appreciate being able to tap into that.

It takes a community of us to keep this alive. Laypeople, lay supporters, monastics, doing what we can do to keep this alive because all the material things, all the ways of the world, those are all so fleeting. If we want to have a happy, peaceful mind when it comes time for us to die and move on, then this is what will carry us through those periods.

This is the long-lasting happiness that we can experience. It's done through Sangha, done through our assembly of noble beings. As the chanting goes, ``They give occasion for incomparable goodness to arise in the world.'' Here we all are, keeping it alive. If we feel despair at times, then it's good to reflect on what it is that we can do to help keep this a sane place to live.
